\input{run_sheets/preamble}
\newcommand{\sheetlabel}{Part 1}
\begin{document}\small
\heading{Part 1: Setting up a project}{Facilitator run sheet}
\textbf{Budget:} 35 min presentation, 10 min Q\&A, 45 min hands-on, 15 min regroup. The hands-on block is the point of the session. Stay out of the way of it.

\topic{Where to stop talking}
\begin{tabularx}{\linewidth}{@{}p{1.45in}X@{}}
\toprule
\textbf{Frame} & \textbf{What to do}\\
\midrule
30 Activity 1 intro & Live demo (below), then stop talking.\\
31 Starter prompt & Stop talking. The 45-minute work block starts.\\
32 Checkpoints & Leave it up while people work.\\
33 Regroup & Participants speak first.\\
\bottomrule
\end{tabularx}

\topic{Before you start}
\begin{itemize}
  \item Have the ZIP link (frame 30) and the starter prompt (frame 31) ready to paste.
  \item Have a spare unzipped copy of the notebook folder open in Claude for the demo.
\end{itemize}

\topic{Lecture, frames 1--29 (35 min)}
The frames carry their own content. If you are behind, shorten frames 28--29 (the reactor case) first: the activity no longer depends on the reactor science.

\topic{Live demo, frames 30--31 (suggested 5 to 7 min)}
\begin{enumerate}
  \item Unzip the download and open the folder itself in Claude.
  \item Paste the starter prompt. Change the wording if you like.
  \item Let Claude inspect the folder. Read its plan aloud.
  \item Ask it a beginner question live, for example: what is a \code{.gitignore}?
  \item Stop it before the first commit. Ask what it would commit, and why.
\end{enumerate}
The point: Claude is the tutor, not a command generator. Do not present a command sequence, and do not explain Git yourself. Say once that the goal is understanding each step, not producing a particular result.

\topic{Work block, 45 min: walk the room}
\textit{Look for these (from the guide's checkpoints):}
\begin{itemize}
  \item The workspace is the project folder itself, not one file or a parent folder.
  \item Claude explains before it acts, and the participant reads the explanation.
  \item The repository is not inside another repository.
  \item The first commit holds the original files as received.
\end{itemize}
\textit{Seen in agent tests, not yet in the room:}
\begin{itemize}
  \item A cautious model paused for confirmation at every step. Tell it how much to do at once: ``Do the next three steps, then stop and summarize.''
  \item A weaker model committed the original files together with its own additions, so there was no untouched starting point.
\end{itemize}
Stuck on a tool for five minutes: pair people up. Pushing to GitHub is optional. A local repository with the push pending is an acceptable checkpoint.

\topic{Regroup, frame 33 (15 min)}
Participants speak first. The questions are on the slide: how far did you get, what did you ask Claude that helped, what surprised you, what will you do next in your research. Ask for a question they would not have known to search for. Note any repository mishaps, so Part 2's recap can address them. Close with ask, inspect, verify, iterate (frame 34).

\vfill
\textbf{After the session:} nothing is required between sessions. The ``keep going'' paths in the Activity 1 guide are optional.
\end{document}
